\subsection*{Parte 5, Misión «Di hola»: retos}
La pregunta de la edad. \py|str( )| vuelve a convertir el número en texto para poder unirlo:
\begin{codenonum}
nombre = input("¿Cómo te llamas? ")
edad = int(input("¿Cuántos años tienes? "))
print("En 10 años tendrás " + str(edad + 10))
\end{codenonum}
\textbf{Súper reto:} sin \py|int( )|, \py|edad + 10| intenta sumar un número a un texto. La última línea dice \texttt{TypeError: can only concatenate str (not "int") to str}. Un \texttt{TypeError} quiere decir que se mezclaron textos y números.

\subsection*{Parte 5, Misión «Adivina mi número»: solo 5 intentos}
\begin{codenonum}
import random
secreto = random.randint(1, 10)
intento = 0
veces = 0
while intento != secreto and veces < 5:
    intento = int(input("Adivina del 1 al 10: "))
    veces += 1
    if intento < secreto:
        print("Muy bajo")
    elif intento > secreto:
        print("Muy alto")
if intento == secreto:
    print("¡Lo adivinaste en " + str(veces) + " intentos!")
else:
    print("¡Se acabaron los intentos! Era " + str(secreto))
\end{codenonum}
Si escribes \texttt{gato} sale \texttt{ValueError: invalid literal for int() with base 10: 'gato'}, porque \py|int( )| no puede convertir la palabra \texttt{gato} en un número.

\subsection*{Parte 5, Misión «Concurso con puntos»: el mensaje final}
Pon esto al final de todo, después de escribir los puntos:
\begin{codenonum}
if puntos == 3:
    print("¡Campeón del concurso!")
elif puntos >= 1:
    print("¡Bien hecho!")
else:
    print("¡Más suerte la próxima vez!")
\end{codenonum}
\textbf{Súper reto:} pon \py|import time| e \py|inicio = time.time()| al principio, y \py|print(round(time.time() - inicio))| al final.

\subsection*{Parte 5, Misión «Prueba de reflejos»: el mejor de tres}
\begin{codenonum}
import time
import random
mejor = 100
for i in range(3):
    input("Presiona Enter para empezar ")
    print("Prepárate...")
    time.sleep(random.randint(2, 5))
    inicio = time.time()
    input("¡YA! ¡Presiona Enter! ")
    tiempo = time.time() - inicio
    if tiempo < mejor:
        mejor = tiempo
print("Tu mejor tiempo: " + str(round(mejor, 2)) + " segundos")
\end{codenonum}
\py|mejor| empieza en 100 segundos, así que tu primer intento seguro lo mejora.

\subsection*{Parte 5: Tradúcelo tú}
Tu código puede verse un poco diferente y aun así estar bien. Si hace lo mismo, ¡lo lograste!
\begin{codenonum}
respuesta = input("¿Cuál es la contraseña? ")
while respuesta != "banana":
    print("¡Mal! Intenta otra vez.")
    respuesta = input("¿Cuál es la contraseña? ")
print("¡Adelante!")
\end{codenonum}
Usar \texttt{while not respuesta == "banana":} en la línea 2 también está bien.

\subsection*{Parte 5: Acertijo de dos líneas}
\begin{codenonum}
edad = int(input("¿Cuántos años tienes? "))
print("El próximo año tendrás " + str(edad + 1))
\end{codenonum}
